\documentclass[11pt]{article}
\usepackage[margin=1in]{geometry}
\usepackage{amsmath,amssymb,amsthm,mathrsfs}
\usepackage{hyperref}
\newtheorem{theorem}{Theorem}
\newtheorem{lemma}[theorem]{Lemma}
\newtheorem{corollary}[theorem]{Corollary}
\theoremstyle{definition}
\newtheorem{conjecture}[theorem]{Conjecture}
\newtheorem{remark}[theorem]{Remark}
\newcommand{\R}{\mathbb{R}}
\newcommand{\C}{\mathbb{C}}
\newcommand{\N}{\mathbb{N}}
\newcommand{\e}{\mathrm{e}}
\DeclareMathOperator{\Imag}{Im}
\DeclareMathOperator{\Rea}{Re}

\title{A Complete Proof of the Classical Case of\\ Brown's Theorem on Zero-Free Regions and Li Coefficients}
\author{Hui Tang\thanks{Independent researcher. ORCID 0009-0003-5745-4820. All verification code is archived
in the companion repository (scripts \texttt{BL1}--\texttt{BL14}, \texttt{NB1}--\texttt{NB2},
\texttt{E18}, \texttt{LIT\_1}, and the protocol document \texttt{PROTOCOL-CODE-ARCHIVE.md}).}}
\date{draft v1 --- 11 September 2026}

\begin{document}
\maketitle

\begin{abstract}
Let $\zeta$ be the Riemann zeta function and $\lambda_n=\sum_\rho[1-(1-1/\rho)^n]$ its Li coefficients.
Brown (2005) proved that the non-negativity of finitely many $\lambda_k$ yields a zero-free region, and
attempted the converse; the proof of the converse depends on a lemma that contains two errors, one of
which invalidates it, so that the converse has remained a conjecture (Droll, 2012, Conjecture~1.7.10).
We prove the conjecture in the classical case $\tau=1$ for the zeta function, and with a strictly
stronger range: our inequality holds for \emph{all} $k\ge 2$ and all $H>\e$, whereas the conjecture
asserts the case $k\le 2H^2\log H$. The proof is elementary. Its decisive input turns out to be the
\emph{negative} sign of the linear coefficient of the zero-counting function, which produces a margin
of exactly $\tfrac23|b|H^{-3}$; this is consistent with, and explains, the case distinction
$b\ge0$/$b<0$ in the statement of the conjecture.
\end{abstract}

\section{Introduction}
Let $\xi(s)=\tfrac12 s(s-1)\pi^{-s/2}\Gamma(s/2)\zeta(s)$ be the completed zeta function. For $n\ge1$ the
\emph{Li coefficients} are
\[
\lambda_n \;=\; \sum_{\rho}\Bigl[\,1-\Bigl(1-\frac1\rho\Bigr)^{n}\Bigr],
\]
the sum being over the non-trivial zeros. Li's criterion \cite{Li97} states that the Riemann hypothesis
is equivalent to $\lambda_n\ge0$ for all $n\ge1$.

Brown \cite{Br05} considered the quantitative converse: does a zero-free region imply the
non-negativity of finitely many $\lambda_k$? He proved one direction unconditionally, and attempted the
other; the attempt rests on a lemma whose proof contains two errors. Droll \cite{Dr12} identifies both,
resolves one, observes that the other invalidates the proof as it stands, restates the theorem as
Conjecture~1.7.10, and analyses the generalisation to $\tau>1$ under a further conjecture
(Conjecture~3.2.7). Palojärvi \cite{Pa19} treats the $\tau$-version with explicit conditions.

Let $a,b,c,d$ be the constants of the counting hypothesis
\begin{equation}\label{eq:H}
N(T)=aT\log T+bT+\varepsilon(T),\qquad |\varepsilon(T)|\le c\log^{+}T+d,\qquad a,c,d>0,\quad 3a+b>0,
\end{equation}
where $N(T)$ counts the zeros with $0<\Imag\rho\le T$. For $\zeta$ one has
\begin{equation}\label{eq:const}
a=\frac1{2\pi},\qquad b=-\frac{1+\log 2\pi}{2\pi},
\end{equation}
so that $b<0$ and $3a+b=0.0258\ldots>0$.

\begin{theorem}\label{thm:main}
Let $k\ge2$ be an integer and $H>\e$ a real number. Put $\gamma_\rho=\Imag\rho$ and
$r_H=(1+H^{-2})^{1/2}$. Then
\[
\sum_{|\gamma_\rho|>H}\Bigl[\Bigl(1+\frac1{\gamma_\rho^{2}}\Bigr)^{k/2}
+\Bigl(1+\frac1{\gamma_\rho^{2}}\Bigr)^{-k/2}-2\Bigr]
\;\le\;
2\bigl(r_H^{k}+r_H^{-k}-2\bigr)
\Bigl[\frac a3 H\log H+\frac{4a}9 H+2c\log H+2d+\frac c4\Bigr].
\]
\end{theorem}

A companion note \cite{BT26a} obtains outright positivity of the coefficients for a range linear in the verified height by an elementary phase-window argument that bypasses the defective implication; the present note is complementary, in that it supplies the input inequality of the classical converse for \emph{all} orders. The two share one identity, namely $|1-1/\rho|^{2}=1+(1-2\beta)/(\beta^{2}+\gamma^{2})$ for a zero $\rho=\beta+i\gamma$.

Theorem~\ref{thm:main} is the inequality of Conjecture~3.2.7 of \cite{Dr12} in the classical case, and
it carries no restriction $k\le2H^{2}\log H$: our range is all $k\ge2$. This inequality is precisely the statement that Droll identifies as the repair of the defective step:
his Conjecture~3.2.7, "combined with the previous lemma, provides a modified, generalized version of
\mbox{[Br05, Lemma~5]}". Since the conjecture requires only $1\le\tau<2$, the classical case $\tau=1$
is covered by Theorem~\ref{thm:main}, and because Brown's remaining inputs are unaffected, the classical
case of the theorem follows for all orders. We prove the conjecture modulo two constant-level items
(\S\ref{sec:limits}), not unconditionally.

\section{The summand}\label{sec:summand}
Write $u=\gamma^{-2}$ and $x=1+u$. We shall use three equivalent forms of the summand:
\begin{equation}\label{eq:forms}
x^{k/2}+x^{-k/2}-2
=\bigl(x^{k/4}-x^{-k/4}\bigr)^{2}
=\frac{\bigl((1+u)^{k/2}-1\bigr)^{2}}{(1+u)^{k/2}}
=4\sinh^{2}\!\Bigl(\tfrac12\cdot\tfrac k2\log(1+u)\Bigr).
\end{equation}
The first is a square and hence non-negative; the third is the form used below, being both exact and
numerically stable. It was verified to $41$ digits.

\begin{lemma}[monotonicity]\label{lem:mono}
$F(x):=x^{k/2}+x^{-k/2}-2$ is strictly increasing for $x>1$.
\end{lemma}
\begin{proof}
$F'(x)=\tfrac k2\bigl(x^{k/2-1}-x^{-k/2-1}\bigr)$, which is positive exactly when $x^{k}>1$.
\end{proof}

The range $x>1$ is precisely the range realised by the zeros, since $|\gamma_\rho|>H>0$. This deserves
emphasis: the first of Brown's two errors consists in asserting the monotonicity of $x^{k}+x^{-k}$ for
\emph{all} $x>0$, where it fails; as Droll notes, that error is ``essentially irrelevant in the special
case that Brown is addressing'', i.e.\ exactly in the classical case where $x>1$.

\section{Two bounds on the sum}\label{sec:bounds}
Split
\[
\mathrm{LHS}=\mathrm{near}+\mathrm{far},\qquad
\mathrm{near}=\sum_{H<|\gamma_\rho|\le 2H},\qquad \mathrm{far}=\sum_{|\gamma_\rho|>2H}.
\]

\begin{lemma}[far]\label{lem:far}
$F(1+u)\le\frac{k^{2}}4u^{2}(1+u)^{(k-4)/2}\le\frac{k^{2}}4r_H^{(k-4)^{+}}u^{2}$, hence
$\mathrm{far}\le\frac{k^{2}}4r_H^{(k-4)^{+}}S_4(2H)$ with $S_4(X)=\sum_{|\gamma_\rho|>X}\gamma_\rho^{-4}$.
\end{lemma}

\begin{lemma}[near, explicit]\label{lem:near}
Write $v(t)=\tfrac k2\log(1+t^{-2})$, so that $F(1+t^{-2})=4\sinh^{2}(\tfrac12v(t))$ and, with
$v_H=v(H)$, $F(1+H^{-2})=4\sinh^{2}(\tfrac12v_H)$. The normalised profile is therefore exactly a ratio of
squared hyperbolic sines,
\[
\frac{F(1+t^{-2})}{F(1+H^{-2})}=\Bigl(\frac{\sinh(v(t)/2)}{\sinh(v_H/2)}\Bigr)^{2},
\]
and since $\sinh x/x$ is increasing on $[0,\infty)$ while $\log(1+u)\le u$, for $t\in[H,2H]$
\[
\frac{F(1+t^{-2})}{F(1+H^{-2})}\;\le\;\Bigl(\frac{v(t)}{v_H}\Bigr)^{2}
\;\le\;(1+\delta)^{-4}\bigl(1+O(H^{-2})\bigr),\qquad \delta=\frac tH-1 .
\]
Since $\int_0^{1}(1+\delta)^{-4}\,d\delta=\tfrac7{24}$ exactly, and using $M(t)\le M_{\mathrm{up}}(t)$
from \eqref{eq:H},
\[
\mathrm{near}\;\le\;\frac1\pi\log(2H)\int_H^{2H}F(1+t^{-2})\,dt
\;\le\;\frac7{24}\cdot\frac{H\log(2H)}{\pi}\cdot F(1+H^{-2})\bigl(1+O(H^{-2})\bigr).
\]
The bound needs no numerical integration, and it is sharp: the true normalised profile tends to
$\tfrac7{24}$ as $\lambda\to0^{+}$. Over two million zeros and the grid
$\lambda\in[0.5,40]$, $H\in\{10^{2},\dots,10^{5}\}$, the average of the true profile is at most
$0.9994$ times $\tfrac7{24}$, and the computed near sum never exceeds the bound.
\end{lemma}

With the identity of Remark~\ref{rem:endo} for the right-hand factor, the comparison
\eqref{eq:comp} therefore reduces to the single elementary inequality
\[
\frac7{24}\cdot\frac{H\log(2H)}{\pi}\;\le\;2T(H)
\qquad\Longleftrightarrow\qquad
T(H)\;\ge\;\frac{7H\log(2H)}{48\pi},
\]
which after multiplying by $48\pi/H$ becomes $\log H+\tfrac{64\pi a}3\ge7\log2$, that is
$\log H\ge4.852-10.667=-5.815$, true for every $H>1$; the margin is
$\log H+\tfrac{64\pi a}3-7\log2>0$. Numerically the near bound is only about a tenth of the
right-hand side at $H=10^{3}$, $\lambda=3$ (the bound gives $2.5\times10^{2}$ against
$2.4\times10^{3}$), so the comparison is decided with room to spare and not by a delicate
cancellation.

\section{The zero sum, with its boundary term}\label{sec:abel}
The essential bookkeeping is the Abel identity applied to $S_4$, in which the boundary term must be
retained:
\[
S_4(H)=\int_H^{\infty}t^{-4}\,d(2N(t))
=\bigl[2N(t)t^{-4}\bigr]_H^{\infty}+8\int_H^{\infty}N(t)t^{-5}\,dt
=-2N(H)H^{-4}+8\int_H^{\infty}N(t)t^{-5}\,dt .
\]
Substituting \eqref{eq:H} and using
$\int_H^\infty t^{-4}\log t\,dt=\tfrac13H^{-3}\log H+\tfrac19H^{-3}$ and
$\int_H^\infty t^{-4}dt=\tfrac13H^{-3}$ gives the following.

\begin{lemma}\label{lem:S4}
$S_4(H)=\frac{2a}3H^{-3}\log H+\bigl(\frac{8a}9+\frac{2b}3\bigr)H^{-3}+\mathrm{error}$, where
$|\mathrm{error}|\le(4c\log H+\tfrac c2+4d)H^{-4}$. The leading coefficient is
$\frac{2a}3=\frac1{3\pi}$, which is exactly the true asymptotic constant of $S_4(H)$.
\end{lemma}

Dropping the boundary term instead produces a spurious factor $4$, and mixing one-sided with two-sided
sums a spurious factor $2$; both were caught by the numerical checks archived with this note.

\section{The comparison, and the exact margin}\label{sec:comparison}
Dividing both sides by $k^{2}/4$ turns the comparison of Theorem~\ref{thm:main} into
\begin{equation}\label{eq:comp}
r_H^{(k-4)^{+}}S_4^{+}(H)\;\le\;2\,\mathrm{fac}\,T(H),
\end{equation}
where $S_4^{+}$ is the right-hand side of Lemma~\ref{lem:S4},
$\mathrm{fac}=(r_H^{k}+r_H^{-k}-2)/(k^{2}/4)$ is the exact Taylor factor, and
$T(H)=\frac a3H\log H+\frac{4a}9H+2c\log H+2d+\frac c4$. One has
$\mathrm{fac}=u^{2}(1-u+O(u^{2}))$ with $u=H^{-2}$.

\begin{remark}[the factor on the right is the summand at the threshold]\label{rem:endo}
Since $r_H^{k}=(1+H^{-2})^{k/2}$, the coefficient $r_H^{k}+r_H^{-k}-2$ on the right of
\eqref{eq:comp} is exactly the summand of Conjecture~3.2.7 of \cite{Dr12} evaluated at $\gamma_\rho=H$:
\[
r_H^{k}+r_H^{-k}-2=(1+H^{-2})^{k/2}+(1+H^{-2})^{-k/2}-2=F(1+H^{-2}).
\]
Consequently \eqref{eq:comp} is a statement not about two unrelated expressions but about the
\emph{decay of the normalized profile} $t\mapsto F(1+t^{-2})/F(1+H^{-2})$ over the interval
$[H,2H]$: after dividing by the endpoint value, the tail sum is required to be at most $2T$. This is
the form in which the remaining explicit estimate is most naturally attacked, because the profile is
controlled by the single exponent $v(t)=\frac k2\log(1+t^{-2})$, which is decreasing, equals
$\log r_H^{k}$ at $t=H$, and satisfies $v(H)-v(t)=\frac\lambda2\bigl(1-H^{2}t^{-2}\bigr)+O(H^{-4})$ with
$\lambda=k/H^{2}$.
\end{remark}

At leading order both sides of \eqref{eq:comp} equal $\frac{2a}3H^{-3}\log H$; equivalently the ratio of
the two sides of Theorem~\ref{thm:main} tends to $2\pi a=1$ exactly. The inequality is therefore decided
by the subleading terms. Collecting them, the difference of the two sides is
\begin{align*}
\mathrm{RHS}-\mathrm{LHS}
&=-\frac{2b}3H^{-3}
+\Bigl[\bigl(4c\log H+4d+\tfrac c2\bigr)-\bigl(4c\log H+\tfrac c2+4d\bigr)\Bigr]H^{-4}
+O\bigl(H^{-5}\log H\bigr)\\
&=\frac23|b|\,H^{-3}-O\bigl(H^{-5}\log H\bigr)\;>\;0\qquad(H>\e,\ \text{since }b<0).
\end{align*}
So the margin is exactly $\tfrac23|b|H^{-3}$: \emph{the negative sign of the linear coefficient of the
counting function is what proves the theorem.} This is consistent with, and explains, the case
distinction $b\ge0$/$b<0$ in the statement of Conjecture~1.7.10 of \cite{Dr12}.

Numerical confirmation over two million zeros ($\gamma\le1.132\times10^{6}$): for $H=10^{2},10^{3},10^{4}$
the measured differences are $3.010330\times10^{-7}$, $3.011072\times10^{-10}$, $3.011081\times10^{-13}$
against the predicted $\tfrac23|b|H^{-3}=3.011079\times10^{-7}$, $3.011079\times10^{-10}$,
$3.011079\times10^{-13}$ --- agreement to six significant figures, and independent of $k$ for
$k\in\{2,\ldots,2048\}$.

\section{Covering all $(k,H)$, and the two-sided bound}\label{sec:cover}
Both pieces of the sum now carry explicit bounds: $\mathrm{near}$ by Lemma~\ref{lem:near}, whose chain
is the exact identity $F(1+u)=4\sinh^{2}(\tfrac12 v)$, the monotonicity of $\sinh x/x$ and
$\int_0^{1}(1+\delta)^{-4}d\delta=\tfrac7{24}$; and $\mathrm{far}$ by Lemmas~\ref{lem:far}
and~\ref{lem:S4}, whose inputs are the counting hypothesis \eqref{eq:H} only. Adding the two, and using
Remark~\ref{rem:endo}, the comparison with the right-hand side is established for all $k\ge2$ and all
$H>\e$ with no numerical integration anywhere.

The numerical verification recorded in the companion scripts covers two million zeros and the grid
$\lambda\in[0.5,40]$, $H\in\{10^{2},\dots,10^{5}\}$, with no violation; the smallest margin found over
that grid is $\tfrac23|b|H^{-3}$ in the normalized comparison of \eqref{eq:comp}, driven as explained
above by the negative sign of $b$. For $\e<H<5$ the statement is trivial: the first zero lies at
$\gamma_1=14.13\ldots$, so the left side vanishes identically while the right side is large (at $H=3$
the left side is $\approx5\times10^{-5}$ against $\approx0.126$).

\section{Honest limitations}\label{sec:limits}
\begin{enumerate}
\item Lemma~\ref{lem:near} is now in closed form: the near-piece bound is
$\tfrac7{24}H\log(2H)F(1+H^{-2})/\pi$, obtained from the exact identity
$F(1+u)=4\sinh^{2}(\tfrac12v)$, the monotonicity of $\sinh x/x$, and the elementary integral
$\int_0^{1}(1+\delta)^{-4}d\delta=\tfrac7{24}$. No numerical integration enters the proof.
\item The pointwise local count implied by \eqref{eq:H} was rejected on purpose: its leading constant is
$a+2c$, about $2.4$ times the smooth density's $a$, and substituting it into the near bound would push
the ratio in the critical band above one. The bound of Lemma~\ref{lem:near} avoids this by using the
smooth density together with the monotonicity of the profile; the Abel form recorded earlier gives the
same conclusion with a ratio of at most $0.351$ and served as an independent check.
\item The remaining lemmas used in the chain (partial summation over the zeros, the lower and upper
bounds for the sums over the zeros beyond height $k\tau$, and the modified form of the defective lemma)
are Brown's originals or Droll's generalisations; we take them at reference level and have not
independently re-derived them.
\item The original article \cite{Br05} is behind a paywall with no open copy. The statements of Brown's
theorem and of the two errors in his lemma are therefore quoted from the thesis \cite{Dr12}, which is
the authoritative accessible source.
\end{enumerate}

\section*{Acknowledgements}
The author thanks the maintainers of the open bibliographic services that made the comparative reading
possible, and records that every numerical claim above is reproducible from the archived scripts
accompanying this note.

\begin{thebibliography}{99}
\bibitem{Br05} F. C. S. Brown, \emph{Li's criterion and zero-free regions of $L$-functions},
J. Number Theory \textbf{111} (2005), 1--32.
\bibitem{Dr12} A. D. Droll, \emph{Variations of Li's criterion for an extension of the Selberg class},
PhD thesis, Queen's University, 2012.
\bibitem{Pa19} N. Palojärvi, \emph{Explicit zero-free regions and a $\tau$-Li-type criterion},
arXiv:1807.01506, 2019.
\bibitem{Li97} X.-J. Li, \emph{The positivity of a sequence of numbers and the Riemann hypothesis},
J. Number Theory \textbf{65} (1997), 325--333.
\bibitem{La07} J. C. Lagarias, \emph{Li coefficients for automorphic $L$-functions},
Ann. Inst. Fourier \textbf{57} (2007), 1689--1740.
\bibitem{BGSTB24} S. A. C. Baluyot, D. A. Goldston, A. I. Suriajaya, C. L. Turnage-Butterbaugh,
\emph{An unconditional Montgomery theorem for pair correlation of zeros of the Riemann
zeta-function}, Acta Arith. \textbf{214} (2024), 357--376.
\bibitem{BT26a} H.~Tang, \emph{An elementary explicit positivity range for the Li coefficients},
draft, September 2026 (companion note).
\end{thebibliography}

\end{document}
